\begin{table}[t]
\centering
\begin{threeparttable}
\caption{Standardized Effect Sizes}
\label{tab:sde}
\begin{adjustbox}{max width=\textwidth}
\begin{tabular}{lcccccc}
\toprule
Outcome & $\hat{\beta}$ & SE & SD($Y$) & SDE & SE(SDE) & Classification \\
\midrule
\multicolumn{7}{l}{\textit{Panel A: Pooled (McCrary discontinuity at each kink)}} \\
Density disc.\ at 79 g/km & -0.712 & (0.910) & 1.890 & -0.377 & (0.482) & Large negative \\
Density disc.\ at 101 g/km & 0.278 & (1.023) & 1.890 & 0.147 & (0.541) & Moderate positive \\
Density disc.\ at 141 g/km & -0.078 & (0.132) & 1.890 & -0.041 & (0.070) & Small negative \\
Density disc.\ at 157 g/km & -0.141 & (0.190) & 1.890 & -0.074 & (0.100) & Moderate negative \\
\addlinespace
\multicolumn{7}{l}{\textit{Panel B: Heterogeneous (Difference-in-Bunching, NL vs.\ DE)}} \\
DiB at 141 (NL$-$DE) & 0.075 & (0.153) & 1.890 & 0.040 & (0.081) & Small positive \\
DiB at 157 (NL$-$DE) & $-$0.015 & (0.209) & 1.890 & $-$0.008 & (0.111) & Small negative \\
\bottomrule
\end{tabular}
\end{adjustbox}
\begin{tablenotes}
\item \textit{Notes:} \textbf{Country:} Netherlands (with Germany as placebo). \textbf{Research question:} Does the Dutch BPM vehicle purchase tax, which features four kinks where marginal CO\textsubscript{2} tax rates increase by 1.6$\times$ to 39.5$\times$, induce manufacturers to manipulate type-approval CO\textsubscript{2} emissions to bunch below thresholds? \textbf{Policy mechanism:} The BPM is a one-time purchase tax calculated from WLTP CO\textsubscript{2} emissions with piecewise-linear rate bands; crossing a kink increases the per-gram rate, raising the marginal cost of higher emissions for buyers and creating incentives for manufacturers to calibrate type-approval emissions below each kink. \textbf{Outcome definition:} McCrary log-density discontinuity at each BPM kink, measuring whether the CO\textsubscript{2} distribution exhibits a break at the threshold. \textbf{Treatment:} Four kink points (79, 101, 141, 157 g/km) with varying marginal rate jumps. \textbf{Data:} European Environment Agency vehicle-level CO\textsubscript{2} monitoring data (Regulation EU 2019/631), 2020--2022, 1.25 million Dutch and 11.5 million German registrations. \textbf{Method:} McCrary density discontinuity test with quadratic in running variable and slope interaction, difference-in-bunching (NL minus DE). \textbf{Sample:} Vehicles with WLTP CO\textsubscript{2} between 1 and 250 g/km; battery-electric vehicles (0 g/km) excluded. SDE $= \hat{\beta} / \text{SD}(Y)$ where SD($Y$) is the standard deviation of log bin counts across all CO\textsubscript{2} values. Classification refers to magnitude, not statistical significance: Large ($|$SDE$| > 0.15$), Moderate ($0.05$--$0.15$), Small ($0.005$--$0.05$), Null ($< 0.005$).
\end{tablenotes}
\end{threeparttable}
\end{table}

